% Generated draft: requires mathematical and source review.
\documentclass[11pt]{article}
\usepackage[margin=25mm]{geometry}
\usepackage{fontspec}
\setmainfont{Latin Modern Roman}
\usepackage{amsmath,amssymb,mathtools,unicode-math}
\setmathfont{Latin Modern Math}
\usepackage{enumitem,needspace,xurl,hyperref}
\setlength{\parindent}{0pt}
\setlength{\parskip}{5pt}
\setlength{\emergencystretch}{3em}
\newcommand{\N}{\mathbb{N}}
\newcommand{\Z}{\mathbb{Z}}
\newcommand{\Q}{\mathbb{Q}}
\newcommand{\R}{\mathbb{R}}
\newcommand{\C}{\mathbb{C}}
\newcommand{\F}{\mathbb{F}}
\newcommand{\A}{\mathbb{A}}
\newcommand{\PP}{\mathbb{P}}
\newcommand{\E}{\mathbb{E}}
\newcommand{\eps}{\varepsilon}
\newcommand{\dd}{\,\mathrm{d}}
\begin{document}
\section*{UnsolvedMath ID: 30002269}
\textbf{Source problem code:} catalogue entry
\subsection{Definitions and mathematical statement}
For the provisional interpretation below, let \(G=(V(G),E(G))\) be a nonempty finite simple undirected graph. A proper vertex coloring assigns different colors to adjacent vertices, and \(\chi(G)\) is the minimum number of colors in such a coloring. For \(v\in V(G)\), define
\[
N_G(v)=\{u\in V(G):\{u,v\}\in E(G)\},\qquad d_G(v)=|N_G(v)|.
\]
Interpret \(\widetilde N(v)\) provisionally as the closed neighborhood \(N_G(v)\cup\{v\}\). For \(S\subseteq V(G)\), let \(G[S]\) denote the induced subgraph on \(S\), and let \(\omega(H)\) denote the maximum number of pairwise adjacent vertices in a graph \(H\).

Under these explicitly provisional conventions, the supplied question is whether
\[
\exists\,\varepsilon\in\R,\ \varepsilon>0,\quad
\forall\,G\text{ nonempty, finite, simple, and undirected},\qquad
\chi(G)\leq
\max_{v\in V(G)}
\left((1-\varepsilon)(d_G(v)+1)
+\varepsilon\,\omega\bigl(G[\widetilde N(v)]\bigr)\right).
\]
The constant \(\varepsilon\) must be independent of \(G\). The displayed bound is retained exactly as supplied, without a ceiling or other rounding operation and without an additional upper restriction on \(\varepsilon\).

The supplied record says ``every graph'' but does not explicitly specify finiteness, simplicity, directedness, exclusion of the empty graph, or the meaning of \(\widetilde N(v)\). Thus the conventions above give a conditional formalization, not a verified identification of the intended domain or notation. For the empty graph the displayed maximum has no supplied convention. The intended neighborhood convention, graph domain, and exact unrounded formulation require source review.
\subsection{Short English statement}
The supplied question asks whether one positive constant, independent of the graph, bounds the chromatic number by the largest vertexwise weighted combination of degree plus one and the clique number in the indicated neighborhood. The graph domain and neighborhood notation are unresolved; the finite simple graph and closed neighborhood interpretation is provisional. The bound is reproduced without rounding.
\subsection{Sources}
\begin{itemize}
\item Supplied local UnsolvedMath registry record, permanent ID \(30002269\), title: Local Clique-Weighted Chromatic Bounds; exact source problem code: \texttt{catalogue entry}. Collection: Oberwolfach Reports Open Problems. The record assigns the problem to 2013 and supplies no named proposer.
\item Supplied primary report citation: Graph Theory (2013), Oberwolfach Reports, DOI 10.4171/owr/2013/02, pp.\ 67--128. \url{https://doi.org/10.4171/owr/2013/02}. The supplied background describes these reports as documenting mathematical workshops and their open-problem discussions.
\item A Local Strengthening of Reed's Conjecture for Quasi-Line Graphs. \url{https://arxiv.org/abs/1109.2112}.
\item A Superlocal Version of Reed's Conjecture. \url{https://arxiv.org/abs/1208.5188}.
\item Graph Theory, Oberwolfach Report 02/2013. \url{https://doi.org/10.4171/owr/2013/02}.
\item Supplied literature assessment, checked on \texttt{2026-08-21}: the registry labels the problem partially solved and reports that the local invariant bound is proved for fractional colorings and important classes such as quasi-line graphs, while remaining open for general graphs. This is supplied dated curation derived from the supplied OWR curation, not an independently verified literature assessment. Its applicability to the exact unrounded formula above requires review; time-sensitive status and cited claims require independent verification.
\item The registry labels the statement \texttt{corrected\_verified} and reports that its mathematical content was checked against the official Oberwolfach report, with notation, context, and question form normalized without changing the source-level claim. This is supplied verification metadata; no independent source verification is asserted here.
\end{itemize}
\end{document}
